% ==========================================================================
%  Guia Completo de Estudo — Prova 2 / Lista II
%  Macroeconomia I (PPGEco-UFES, 2026)
%  Prof. Ricardo Ramalhete Moreira
%  Rigidez Nominal (cap.6) e NK DSGE (cap.7)
% ==========================================================================
\documentclass[12pt, a4paper]{article}

% --------------------------------------------------------------------------
% Packages
% --------------------------------------------------------------------------
\usepackage[utf8]{inputenc}
\usepackage[T1]{fontenc}
\usepackage{lmodern}
\usepackage[brazil]{babel}
\usepackage{microtype}

\usepackage[top=2.5cm, bottom=2.5cm, left=3cm, right=3cm, headheight=15pt]{geometry}

\usepackage{amsmath, amssymb, amsthm}
\usepackage{mathtools}
\usepackage{bm}

\usepackage{booktabs}
\usepackage{array}
\usepackage{multicol}
\usepackage{xcolor}
\usepackage{hyperref}
\usepackage{enumitem}
\usepackage{tcolorbox}
\tcbuselibrary{breakable}

% --------------------------------------------------------------------------
% Cores
% --------------------------------------------------------------------------
\definecolor{cyanrbc}{HTML}{3b9dbd}
\definecolor{orangerbc}{HTML}{d97b39}
\definecolor{grayrbc}{HTML}{6b7280}
\definecolor{greenbox}{HTML}{166534}
\definecolor{greenbg}{HTML}{dcfce7}
\definecolor{boxbg}{HTML}{f5f0e8}
\definecolor{boxborder}{HTML}{3b9dbd}
\definecolor{provabg}{HTML}{fff7ed}
\definecolor{provaborder}{HTML}{d97b39}

% --------------------------------------------------------------------------
% Caixas (standard skin — sem skins/TikZ)
% --------------------------------------------------------------------------
\newtcolorbox{resultbox}[1][]{standard, colback=boxbg, colframe=boxborder,
  boxrule=1pt, arc=3pt, left=6pt, right=6pt, top=4pt, bottom=4pt,
  fonttitle=\bfseries\small, #1}
\newtcolorbox{provabox}[1][]{standard, colback=provabg, colframe=provaborder,
  boxrule=1.2pt, arc=3pt, left=6pt, right=6pt, top=4pt, bottom=4pt,
  fonttitle=\bfseries\small, #1}
\newtcolorbox{formulabox}[1][]{standard, colback=greenbg!30, colframe=greenbox!60,
  boxrule=0.8pt, arc=3pt, left=6pt, right=6pt, top=4pt, bottom=4pt,
  fonttitle=\bfseries\small, #1}

% --------------------------------------------------------------------------
% Comandos matemáticos
% --------------------------------------------------------------------------
\newcommand{\E}{\mathbb{E}}
\newcommand{\dif}{\mathrm{d}}
\newcommand{\pibar}{\bar{\pi}}
\newcommand{\xgap}{x_t}

% --------------------------------------------------------------------------
% Espaçamento e parágrafos
% --------------------------------------------------------------------------
\setlength{\parskip}{0.5\baselineskip}
\setlength{\parindent}{0pt}

% --------------------------------------------------------------------------
% Títulos
% --------------------------------------------------------------------------
\usepackage{titlesec}
\titleformat{\section}{\large\bfseries\color{cyanrbc}}{Questão \thesection.}{0.5em}{}[\vspace{-0.3em}\rule{\linewidth}{0.4pt}]
\titleformat{\subsection}{\normalsize\bfseries\color{grayrbc}}{}{0em}{}

% --------------------------------------------------------------------------
% Cabeçalho/rodapé
% --------------------------------------------------------------------------
\usepackage{fancyhdr}
\pagestyle{fancy}
\fancyhf{}
\lhead{\footnotesize\textcolor{grayrbc}{Guia Completo P2 --- Macroeconomia I (PPGEco-UFES, 2026)}}
\rhead{\footnotesize\textcolor{grayrbc}{Rigidez Nominal \(\cdot\) NK DSGE --- Romer, cap.\ 6--7}}
\cfoot{\small\thepage}
\renewcommand{\headrulewidth}{0.4pt}

% ==========================================================================
\begin{document}
% ==========================================================================

% --------------------------------------------------------------------------
% Capa
% --------------------------------------------------------------------------
\begin{titlepage}
\begin{center}
\vspace*{1.5cm}
{\Large\bfseries Universidade Federal do Espírito Santo}\\[0.5em]
{\large Programa de Pós-Graduação em Economia (PPGEco-UFES)}\\[2em]

\rule{\linewidth}{1.2pt}\\[0.8em]
{\LARGE\bfseries\color{cyanrbc} Guia Completo de Estudo --- Prova 2 / Lista II\\[0.4em]
\large Macroeconomia I --- 2026}\\[0.8em]
\rule{\linewidth}{1.2pt}

\vspace{0.8em}
{\large Romer, caps.\ 6 e 7}

\vspace{0.8em}
{\large Prof.\ Ricardo Ramalhete Moreira}

\vspace{1.5em}
\begin{tcolorbox}[standard, colback=provabg, colframe=provaborder, width=0.88\textwidth,
  arc=4pt, boxrule=1.2pt, left=8pt, right=8pt, top=6pt, bottom=6pt]
\centering
\textbf{\textcolor{provaborder}{Cobertura completa:}} Este guia cobre \textbf{TODOS} os
tópicos de P2: \textbf{Rigidez Nominal} (Romer, cap.\ 6) e
\textbf{Modelos DSGE Novo-Keynesianos} (Romer, cap.\ 7). Inclui derivações
passo a passo e soluções completas das 11 questões da Lista II.
\end{tcolorbox}

\vspace{1.5em}
\begin{tabular}{lll}
\toprule
\textbf{Questão} & \textbf{Tópico} & \textbf{Capítulo} \\
\midrule
Q1 & Custos de menu e decisão de ajuste de preços & Cap.\ 6 \\
Q2 & Externalidade de demanda e falha de coordenação & Cap.\ 6 \\
Q3 & Modelo de Calvo: probabilidade e duração & Cap.\ 6 \\
Q4 & Derivação da Curva de Phillips NK & Cap.\ 6 \\
Q5 & Rigidezes reais vs.\ nominais & Cap.\ 6 \\
Q6 & IS dinâmica e taxa natural & Cap.\ 7 \\
Q7 & CPNK: trade-off e choque de custo-push & Cap.\ 7 \\
Q8 & Regra de Taylor e princípio de Taylor & Cap.\ 7 \\
Q9 & Solução NK: choque de demanda & Cap.\ 7 \\
Q10 & Solução NK: choque de custo-push & Cap.\ 7 \\
Q11 & Fronteira eficiente de política monetária & Cap.\ 7 \\
\bottomrule
\end{tabular}

\vspace{1.5em}
\begin{tcolorbox}[standard, colback=boxbg, colframe=boxborder, width=0.88\textwidth,
  arc=4pt, boxrule=1pt, left=8pt, right=8pt, top=6pt, bottom=6pt]
\centering
\textbf{Referência principal:} \\[0.3em]
D.\ Romer, \textit{Advanced Macroeconomics}, 5.ª ed., McGraw-Hill (2019). \\
Capítulos 6 e 7. \\[0.3em]
J.\ Galí, \textit{Monetary Policy, Inflation, and the Business Cycle}, Princeton (2008).
\end{tcolorbox}

\vfill
{\small Soluções computadas com Python; código disponível em
\texttt{github.com/fcarva/macro}.}
\end{center}
\end{titlepage}

% --------------------------------------------------------------------------
% Formulário de Referência Rápida
% --------------------------------------------------------------------------
\newpage
\thispagestyle{fancy}

\begin{center}
{\large\bfseries\color{cyanrbc} Formulário de Referência Rápida --- P2}\\[0.3em]
\rule{\linewidth}{0.6pt}
\end{center}

\vspace{0.3em}
\begin{multicols}{2}

\textbf{\textcolor{cyanrbc}{BLOCO RIGIDEZ NOMINAL (Cap.\ 6)}}

\medskip
\textbf{Demanda Dixit-Stiglitz:}
\[
Y_i = Y\!\left(\frac{P_i}{P}\right)^{-\eta}, \quad \eta > 1
\]

\textbf{Preço ótimo flexível:}
\[
P_i^* = \mu\cdot MC_i,\quad \mu = \frac{\eta}{\eta-1}
\]

\textbf{Ganho de reajuste (Mankiw):}
\[
G = \frac{\eta-1}{2}\,p_{dev}^2
\]

\textbf{Limiar de ajuste:}
\[
|p_{dev}| \geq \bar{p} = \sqrt{\frac{2z}{\eta-1}}
\]

\textbf{Calvo --- preço de re-otimização:}
\[
\log\tilde{P}_t = (1-\alpha\beta)\!\sum_{s=0}^{\infty}(\alpha\beta)^s
  \E_t[\log P_{t+s}^*]
\]

\textbf{Índice de preços (log):}
\[
p_t = \alpha\,p_{t-1} + (1-\alpha)\,\tilde{p}_t
\]

\textbf{CPNK:}
\[
\pi_t = \beta\,\E_t[\pi_{t+1}] + \kappa\,x_t + u_t
\]

\textbf{Slope:}
\[
\kappa = \frac{(1-\alpha)(1-\alpha\beta)}{\alpha}\,(\omega + \sigma^{-1})
\]

\columnbreak

\textbf{\textcolor{cyanrbc}{BLOCO NK DSGE (Cap.\ 7)}}

\medskip
\textbf{IS dinâmica:}
\[
x_t = \E_t[x_{t+1}] - \frac{1}{\sigma}(i_t - \E_t[\pi_{t+1}] - r^n)
\]

\textbf{CPNK:}
\[
\pi_t = \beta\,\E_t[\pi_{t+1}] + \kappa\,x_t + u_t
\]

\textbf{Regra de Taylor:}
\[
i_t = r^n + \phi_\pi\,\pi_t + \phi_x\,x_t + v_t
\]

\textbf{Taxa natural:}
\[
r^n = -\ln\beta \approx \frac{1-\beta}{\beta}
\]

\textbf{Solução choque $v_t$:}
\[
\psi_{xv} = \frac{-1}{\sigma\Delta_v},\quad
\psi_{\pi v} = \frac{-\kappa}{\sigma\Delta_v(1-\beta\rho_v)}
\]

\textbf{Solução choque $u_t$:}
\[
\psi_{xu} = \frac{-(\phi_\pi-\rho_u)}{\sigma(1-\beta\rho_u)D_u}
\]

\textbf{Princípio de Taylor:} $\phi_\pi > 1$

\textbf{Condição precisa:}
\[
\kappa(\phi_\pi-1) + (1-\beta)\phi_x > 0
\]

\textbf{Calibração padrão:}\\[2pt]
$\beta=0{,}99$, $\sigma=1$, $\kappa=0{,}1$, $\phi_\pi=1{,}5$, $\phi_x=0{,}5$

\end{multicols}

\rule{\linewidth}{0.6pt}

\vspace{0.5em}
\begin{provabox}[title={Princípio de Taylor --- Alerta de Prova}]
\textbf{$\phi_\pi > 1$} é condição necessária para determinação do equilíbrio.
Se $\phi_\pi < 1$ (política \emph{passiva}), a taxa real cai quando a inflação sobe,
realimentando as expectativas inflacionárias --- equilíbrio \textbf{indeterminado}
(múltiplos EE). A armadilha de liquidez ocorre quando $i_t = 0$ impede atingir
$r^n < 0$: o BC perde a capacidade de estimular a demanda via juros convencional.
\end{provabox}

% --------------------------------------------------------------------------
% Sumário
% --------------------------------------------------------------------------
\newpage
\tableofcontents
\newpage

% ==========================================================================
%  PARTE 0 — Contexto
% ==========================================================================

\section*{\large\color{cyanrbc} Parte 0 --- Contexto: Por que Rigidez Nominal Importa?}
\addcontentsline{toc}{section}{Parte 0 --- Contexto: Por que Rigidez Nominal Importa?}

Em mercados competitivos com preços plenamente flexíveis, a moeda é neutra: um choque
monetário eleva apenas o nível nominal de preços, deixando todas as variáveis reais
(produto, emprego, salário real) inalteradas. Isso é a \textbf{dicotomia clássica}.

A Macroeconomia Novo-Keynesiana parte de uma premissa diferente: preços são
\textbf{rígidos no curto prazo}. Essa rigidez --- por custos de menu (Mankiw, 1985),
por contratos escalonados (Fischer, 1977; Taylor, 1980) ou por chegadas aleatórias de
oportunidades de ajuste (Calvo, 1983) --- implica que choques monetários têm efeitos
\textbf{reais} de curto prazo.

\textbf{Conexão entre capítulos 6 e 7:} O cap.\ 6 fundamenta \emph{por que} os preços
são rígidos e deriva a Curva de Phillips NK (CPNK). O cap.\ 7 usa a CPNK como bloco de
oferta do modelo de três equações, que se torna o framework central para análise de
política monetária em economias modernas.

\begin{resultbox}[title={Lógica da sequência}]
\begin{enumerate}[leftmargin=1.5em, itemsep=2pt]
  \item Firmas têm poder de mercado (Dixit-Stiglitz) $\Rightarrow$ markup $\mu$.
  \item Custos de menu ou Calvo $\Rightarrow$ preços rígidos no CP.
  \item Rigidez $\Rightarrow$ moeda não-neutra $\Rightarrow$ CPNK com hiato do produto.
  \item CPNK + IS dinâmica + Regra de Taylor $\Rightarrow$ Modelo NK.
  \item Modelo NK $\Rightarrow$ análise de PM ótima, princípio de Taylor, IRFs.
\end{enumerate}
\end{resultbox}

% ==========================================================================
%  PARTE I — Rigidez Nominal (Derivações, Cap. 6)
% ==========================================================================

\newpage
\section*{\large\color{cyanrbc} Parte I --- Rigidez Nominal: Derivações (Cap.\ 6)}
\addcontentsline{toc}{section}{Parte I --- Rigidez Nominal: Derivações (Cap.\ 6)}

% ------------------------------------------------------------------
\subsection*{I.1 \ Competição Monopolística e Precificação Ótima}

\textbf{Demanda Dixit-Stiglitz.}
O agregador CES da economia gera demanda para cada variedade $i$:
\[
Y = \left[\int_0^1 Y_i^{(\eta-1)/\eta}\,di\right]^{\eta/(\eta-1)},\quad \eta > 1.
\]
Maximizando a utilidade do consumidor para um dado gasto $PY$, obtém-se
a demanda pela variedade $i$:
\begin{equation}
Y_i = Y\!\left(\frac{P_i}{P}\right)^{-\eta},
\label{eq:demand_ds}
\end{equation}
onde $P = \left[\int_0^1 P_i^{1-\eta}\,di\right]^{1/(1-\eta)}$ é o índice de preços
Dixit-Stiglitz e $\eta$ é a elasticidade-preço.

\textbf{Preço ótimo — derivação passo a passo.}
O lucro da firma $i$ é:
\[
\Pi_i = (P_i - MC_i)\,Y_i = (P_i - MC_i)\,Y\!\left(\frac{P_i}{P}\right)^{-\eta}.
\]

\textit{Passo 1.} Derivada em relação a $P_i$:
\[
\frac{\partial\Pi_i}{\partial P_i} = Y\!\left(\frac{P_i}{P}\right)^{-\eta}
  + (P_i - MC_i)(-\eta)\,Y\,P_i^{-\eta-1}\,P^\eta = 0.
\]

\textit{Passo 2.} Fatorar $Y(P_i/P)^{-\eta}$:
\[
Y\!\left(\frac{P_i}{P}\right)^{-\eta}\!\!\left[1 - \frac{\eta(P_i - MC_i)}{P_i}\right] = 0.
\]

\textit{Passo 3.} Resolver (o fator entre colchetes deve ser zero):
\[
1 - \frac{\eta(P_i - MC_i)}{P_i} = 0
\;\Rightarrow\;
P_i - MC_i = \frac{P_i}{\eta}
\;\Rightarrow\;
P_i\!\left(1 - \frac{1}{\eta}\right) = MC_i.
\]

\begin{resultbox}[title={Preço ótimo com poder de mercado}]
\[
\boxed{P_i^* = \frac{\eta}{\eta-1}\,MC_i \equiv \mu\cdot MC_i,
\quad \mu = \frac{\eta}{\eta-1} > 1.}
\]
Em log: $p_i^* = \log\mu + mc_i$.
O markup $\mu$ é decrescente em $\eta$: mais concorrência $\Rightarrow$ menor markup.
\end{resultbox}

% ------------------------------------------------------------------
\subsection*{I.2 \ Custos de Menu (Mankiw 1985)}

\textbf{Ganho de reajuste — derivação quadrática.}
Expandindo o lucro em segunda ordem ao redor do preço ótimo $P_i^*$:
\[
\Pi(P_i) \approx \Pi(P_i^*) + \underbrace{\Pi'(P_i^*)}_{=\,0}\,(P_i - P_i^*)
  + \frac{1}{2}\Pi''(P_i^*)\,(P_i - P_i^*)^2.
\]
A derivada de primeira ordem é nula no ótimo. Em termos de log-desvio
$p_{dev} = p_i - p_i^*$, normalizando a derivada de segunda ordem:
\[
\Pi(P_i) \approx \Pi(P_i^*) - \frac{\eta-1}{2}\,p_{dev}^2.
\]
O \textbf{ganho de reajustar} (diferença de lucros) é:

\begin{formulabox}[title={Ganho de reajuste (Mankiw 1985)}]
\[
G(p_{dev}) = \Pi(P_i^*) - \Pi(P_i) = \frac{\eta-1}{2}\,p_{dev}^2.
\]
O ganho é \emph{quadrático} em $p_{dev}$: desvios pequenos têm custo de
segundo ordem para a firma, mas de primeiro ordem para o bem-estar social.
\end{formulabox}

\textbf{Limiar de ajuste.}
A firma ajusta se $G(p_{dev}) \geq z$ (custo de menu):
\[
\frac{\eta-1}{2}\,p_{dev}^2 \geq z
\;\Longleftrightarrow\;
|p_{dev}| \geq \bar{p}(z) = \sqrt{\frac{2z}{\eta-1}}.
\]

\begin{resultbox}[title={Limiar de ajuste de preços}]
\[
\boxed{|p_{dev}| \geq \bar{p}(z) = \sqrt{\frac{2z}{\eta-1}}.}
\]
\begin{itemize}[leftmargin=1.5em, itemsep=2pt]
  \item $\bar{p}(z)$ cresce com $z$: custos maiores $\Rightarrow$ mais rigidez.
  \item $\bar{p}(z)$ decresce com $\eta$: demanda mais elástica $\Rightarrow$
        firmas ajustam com mais frequência.
\end{itemize}
\end{resultbox}

\textbf{Externalidade de demanda agregada (Ball-Romer).}
Com $n$ firmas simétricas, o ganho \emph{social} de ajuste é:
\[
G^{social} = n\cdot G^{private} = n\cdot\frac{\eta-1}{2}\,p_{dev}^2.
\]
\textbf{Região de falha de coordenação:} $G^{private} < z \leq G^{social}$.
Nessa região, nenhuma firma ajusta individualmente (pois $G^{private} < z$),
mas o ajuste coletivo seria benéfico (pois $G^{social} \geq z$).
Resultado: múltiplos equilíbrios e rigidez agregada mesmo com custos individuais ínfimos.

\begin{provabox}[title={Resultado central de Mankiw-Ball-Romer}]
Custos de menu \textbf{pequenos} geram rigidez \textbf{grande} por dois motivos:
(1) a externalidade de demanda amplifica o ganho social; (2) firmas são
\emph{complementos estratégicos} --- se as demais não ajustam, é ótimo não ajustar
também (o preço relativo da firma desviante muda pouco).
\end{provabox}

% ------------------------------------------------------------------
\subsection*{I.3 \ Modelo de Calvo}

\textbf{Estrutura.}
Em cada período, cada firma:
\begin{itemize}[leftmargin=1.5em, itemsep=2pt]
  \item mantém o preço com probabilidade $\alpha \in (0,1)$, ou
  \item re-otimiza com probabilidade $(1-\alpha)$.
\end{itemize}
O processo é i.i.d.: independente do estado e do tempo desde o último ajuste.

A probabilidade de que um preço fixado há $s$ períodos ainda vigore é $\alpha^s$.
A duração esperada de um preço é:
\[
E[T] = \sum_{s=0}^{\infty}(1-\alpha)\alpha^s \cdot s = \frac{\alpha}{1-\alpha}
\;\Rightarrow\; E[T]+1 = \frac{1}{1-\alpha}.
\]
Para $\alpha = 0{,}75$ (trimestral): $E[T] = 3$ períodos adicionais, duração total
= 4 trimestres = 1 ano.

\textbf{Preço de re-otimização.}
A firma que re-otimiza em $t$ escolhe $\tilde{P}_t$ para maximizar o valor presente
esperado dos lucros, ponderado pela probabilidade $\alpha^s$ de o preço ainda viger:
\begin{equation}
\log\tilde{P}_t = (1-\alpha\beta)\sum_{s=0}^{\infty}(\alpha\beta)^s
  \E_t\!\left[\log P_{t+s}^*\right].
\label{eq:calvo_reset}
\end{equation}
O peso $(1-\alpha\beta)$ garante que $\sum_{s=0}^\infty (\alpha\beta)^s(1-\alpha\beta) = 1$
(normalização).

\begin{formulabox}[title={Preço de re-otimização de Calvo}]
$\tilde{P}_t$ é uma \textbf{média ponderada} dos preços ótimos futuros esperados,
com pesos que decrescem geometricamente à taxa $\alpha\beta$.
\end{formulabox}

\textbf{Índice de preços --- log-linearização.}
O índice CES agrega preços antigos (peso $\alpha$) e novos (peso $1-\alpha$):
\[
P_t^{1-\eta} = \alpha\,P_{t-1}^{1-\eta} + (1-\alpha)\,\tilde{P}_t^{1-\eta}.
\]
Log-linearizando em torno de inflação zero:
\begin{equation}
p_t = \alpha\,p_{t-1} + (1-\alpha)\,\tilde{p}_t.
\label{eq:price_index}
\end{equation}
Equivalentemente: $\pi_t = p_t - p_{t-1} = (1-\alpha)(\tilde{p}_t - p_{t-1})$,
ou seja, a inflação é a fração de firmas que ajusta multiplicada pelo gap de ajuste.

% ------------------------------------------------------------------
\subsection*{I.4 \ Derivação da CPNK}

A seguir apresentamos a derivação completa da Curva de Phillips NK a partir da
precificação de Calvo.

\textbf{Passo 1.} Custo marginal em desvio log do estado estacionário:
$\hat{mc}_t$. O preço ótimo flexível é $\log P_t^* = \log\mu + p_t + \hat{mc}_t$,
onde $p_t$ é o nível de preços e $\hat{mc}_t$ captura custos reais.

\textbf{Passo 2.} Substituindo em \eqref{eq:calvo_reset}:
\[
\tilde{p}_t - p_t = (1-\alpha\beta)\sum_{s=0}^{\infty}(\alpha\beta)^s
  \E_t\!\left[\hat{mc}_{t+s} + (p_{t+s} - p_t)\right].
\]

\textbf{Passo 3.} Note que $p_{t+s} - p_t = \sum_{j=1}^s \pi_{t+j}$.

\textbf{Passo 4.} Definir o \emph{reset price gap}: $\hat{q}_t = \tilde{p}_t - p_{t-1}$.
Da equação do índice \eqref{eq:price_index}:
\[
\pi_t = (1-\alpha)\,\hat{q}_t.
\]

\textbf{Passo 5.} Manipulação recursiva. O preço de re-otimização satisfaz a recursão:
\[
\hat{q}_t = (1-\alpha\beta)\,\hat{mc}_t + \alpha\beta\,\E_t[\hat{q}_{t+1}]
  + \alpha\beta\,\E_t[\pi_{t+1}] + \pi_t.
\]
Usando $\pi_t = (1-\alpha)\hat{q}_t$ e $\pi_{t+1} = (1-\alpha)\hat{q}_{t+1}$:
\[
\frac{\pi_t}{1-\alpha} = (1-\alpha\beta)\,\hat{mc}_t
  + \alpha\beta\,\E_t\!\left[\frac{\pi_{t+1}}{1-\alpha}\right]
  + \alpha\beta\,\E_t[\pi_{t+1}] + \pi_t.
\]

\textbf{Passo 6.} Reorganizando (dividindo por $1-\alpha$ e simplificando termos):
\[
\pi_t\!\left(\frac{1}{1-\alpha} - 1\right) = (1-\alpha\beta)\,\hat{mc}_t
  + \frac{\alpha\beta}{1-\alpha}\,\E_t[\pi_{t+1}]
  + \alpha\beta\,\E_t[\pi_{t+1}].
\]
\[
\frac{\alpha}{1-\alpha}\,\pi_t = (1-\alpha\beta)\,\hat{mc}_t
  + \frac{\alpha\beta(2-\alpha)}{1-\alpha}\,\E_t[\pi_{t+1}].
\]
Multiplicando por $(1-\alpha)/\alpha$:
\[
\pi_t = \frac{(1-\alpha)(1-\alpha\beta)}{\alpha}\,\hat{mc}_t + \beta\,\E_t[\pi_{t+1}].
\]

\textbf{Passo 7.} Usando $\hat{mc}_t = (\omega + \sigma^{-1})x_t$ (do lado real do modelo):

\begin{formulabox}[title={Curva de Phillips NK (CPNK)}]
\[
\boxed{\pi_t = \beta\,\E_t[\pi_{t+1}] + \kappa\,x_t + u_t,}
\]
\[
\kappa = \frac{(1-\alpha)(1-\alpha\beta)}{\alpha}\,(\omega + \sigma^{-1}).
\]
\textbf{Interpretação de $\kappa$:}
\begin{itemize}[leftmargin=1.5em, itemsep=2pt]
  \item $\alpha\nearrow$ (mais rigidez) $\Rightarrow$ $\kappa\searrow$ (CPNK mais plana).
  \item $\beta\nearrow$ (firmas mais pacientes) $\Rightarrow$ expectativas mais importantes.
  \item $\omega\nearrow$ ou $\sigma\searrow$ $\Rightarrow$ $\kappa\nearrow$ (pass-through maior).
\end{itemize}
\textbf{Calibração:} $\alpha=0{,}75$, $\beta=0{,}99$, $\omega=1$, $\sigma=1$:
$\kappa = (0{,}25)(0{,}2575)/0{,}75 \times 2 \approx 0{,}085$.
(Galí 2008 usa $\kappa=0{,}1$ com parâmetros ligeiramente diferentes.)
\end{formulabox}

% ------------------------------------------------------------------
\subsection*{I.5 \ Rigidezes Reais e Falha de Coordenação}

\textbf{Rigidezes nominais vs.\ reais.}
\begin{itemize}[leftmargin=1.5em, itemsep=2pt]
  \item \textbf{Rigidez nominal}: custo de alterar o preço nominal (custo de menu $z$).
  \item \textbf{Rigidez real}: preço ótimo relativo $P_i^*/P$ responde pouco a variações
        macroeconômicas. Medida pelo parâmetro $\gamma = \partial P_i^*/\partial P \in (0,1)$
        (complementaridade estratégica).
\end{itemize}

\textbf{Mecanismo de complementação.}
Com $\gamma$ próximo de 1 (forte complementaridade):
\begin{enumerate}[leftmargin=1.5em, itemsep=2pt]
  \item Demanda sobe $\Rightarrow$ $P^*$ deveria subir, mas pouco (rigidez real).
  \item Logo $p_{dev} = p_i - p_i^*$ é pequeno.
  \item $G = (\eta-1)/2 \cdot p_{dev}^2$ é ínfimo.
  \item Qualquer $z > 0$ é suficiente para impedir ajuste.
\end{enumerate}
\textbf{Conclusão de Ball-Romer (1990):} rigidezes reais amplificam os efeitos de
custos de menu pequenos, tornando a rigidez agregada de preços muito mais robusta.

\begin{provabox}[title={Ball-Romer: multiplicação de rigidezes}]
Pense assim: rigidez real $\Rightarrow$ $P^*$ não muda muito $\Rightarrow$
custo de não ajustar é baixo $\Rightarrow$ qualquer custo de menu, por menor que seja,
é suficiente para impedir o ajuste. As duas rigidezes se \emph{multiplicam}:
uma pequena rigidez de cada tipo gera uma grande rigidez agregada.
\end{provabox}

% ==========================================================================
%  PARTE II — Modelo NK Canônico (Derivações, Cap. 7)
% ==========================================================================

\newpage
\section*{\large\color{cyanrbc} Parte II --- Modelo NK Canônico: Derivações (Cap.\ 7)}
\addcontentsline{toc}{section}{Parte II --- Modelo NK Canônico: Derivações (Cap.\ 7)}

% ------------------------------------------------------------------
\subsection*{II.1 \ Três Equações do Modelo NK}

O modelo Novo-Keynesiano canônico consiste em três equações log-linearizadas:

\textbf{IS dinâmica} (equação de Euler das famílias):
\begin{equation}
x_t = \E_t[x_{t+1}] - \frac{1}{\sigma}(i_t - \E_t[\pi_{t+1}] - r^n).
\label{eq:is_nk}
\end{equation}
\textit{Derivação:} A família CRRA maximiza $\sum_t \beta^t u(C_t)$ com
$u'(C_t) = \beta(1+r_{t+1})u'(C_{t+1})$. Log-linearizando ($\hat{c}_t = x_t$ no
modelo simples, sem gasto público):
$x_t = \E_t[x_{t+1}] - (1/\sigma)(i_t - \E_t[\pi_{t+1}] - r^n)$,
onde $r^n = -\ln\beta \approx (1-\beta)/\beta$ é a taxa real natural e $1/\sigma$
é a elasticidade de substituição intertemporal (EIS).

Para $\beta=0{,}99$: $r^n \approx 1\%/\text{trimestre} = 4\%/\text{ano}$.

\textbf{CPNK} (do cap.\ 6):
\begin{equation}
\pi_t = \beta\,\E_t[\pi_{t+1}] + \kappa\,x_t + u_t.
\label{eq:cpnk_nk}
\end{equation}

\textbf{Regra de Taylor}:
\begin{equation}
i_t = r^n + \phi_\pi\,\pi_t + \phi_x\,x_t + v_t, \quad \phi_\pi > 1,\;\phi_x \geq 0.
\label{eq:taylor_nk}
\end{equation}
$v_t$ é o choque de política monetária (distúrbio de demanda);
$u_t$ é o choque de custo-push (distúrbio de oferta).

\begin{resultbox}[title={Interpretação das três equações}]
\begin{enumerate}[leftmargin=1.5em, itemsep=3pt]
  \item \textbf{IS:} Euler intertemporal. $i_t - \E_t[\pi_{t+1}] > r^n$
    (taxa real acima do natural) $\Rightarrow$ famílias poupam mais $\Rightarrow$
    $x_t < \E_t[x_{t+1}]$ (produto abaixo da trajetória esperada).
  \item \textbf{CPNK:} forward-looking. A inflação atual depende da inflação
    \emph{futura} esperada e do hiato do produto corrente.
  \item \textbf{Taylor:} reação do BC. $\phi_\pi > 1$ garante que taxa real suba
    quando $\pi$ sobe --- estabilizador automático.
\end{enumerate}
\end{resultbox}

% ------------------------------------------------------------------
\subsection*{II.2 \ Solução por Coeficientes Indeterminados}

\textbf{Sistema IS modificado.}
Substituindo \eqref{eq:taylor_nk} em \eqref{eq:is_nk}:
\begin{equation}
\left(1 + \frac{\phi_x}{\sigma}\right)x_t + \frac{\phi_\pi}{\sigma}\,\pi_t
= \E_t[x_{t+1}] + \frac{1}{\sigma}\,\E_t[\pi_{t+1}] - \frac{v_t}{\sigma}.
\label{eq:is_mod}
\end{equation}

\textbf{Choque de demanda $v_t = \rho_v v_{t-1} + \varepsilon_t^v$, $u_t=0$.}

\textit{Passo 1.} Adivinhar $x_t = \psi_{xv}v_t$, $\pi_t = \psi_{\pi v}v_t$.

\textit{Passo 2.} Da CPNK \eqref{eq:cpnk_nk} com $u_t=0$:
\[
\psi_{\pi v}v_t = \beta\rho_v\psi_{\pi v}v_t + \kappa\psi_{xv}v_t
\;\Rightarrow\;
\psi_{\pi v}(1-\beta\rho_v) = \kappa\psi_{xv}
\;\Rightarrow\;
\psi_{\pi v} = \frac{\kappa\psi_{xv}}{1-\beta\rho_v}. \quad (*)
\]

\textit{Passo 3.} Da IS modificada \eqref{eq:is_mod}, dividindo por $v_t$
(com $\E_t[v_{t+1}]=\rho_v v_t$):
\[
\psi_{xv}\!\left(1 + \frac{\phi_x}{\sigma} - \rho_v\right)
+ \psi_{\pi v}\!\frac{\phi_\pi - \rho_v}{\sigma} = -\frac{1}{\sigma}.
\]

\textit{Passo 4.} Substituindo $(*)$:
\[
\psi_{xv}\!\left[(1-\rho_v) + \frac{1}{\sigma}\!\left(\phi_x
  + \frac{\kappa(\phi_\pi-\rho_v)}{1-\beta\rho_v}\right)\right] = -\frac{1}{\sigma}.
\]

\textit{Passo 5.} Definindo $\Delta_v = (1-\rho_v) + \frac{1}{\sigma}\!\left[\phi_x + \frac{\kappa(\phi_\pi-\rho_v)}{1-\beta\rho_v}\right]$:

\begin{formulabox}[title={Solução para choque de demanda $v_t$}]
\[
\boxed{
\psi_{xv} = \frac{-1}{\sigma\Delta_v},\qquad
\psi_{\pi v} = \frac{-\kappa}{\sigma\Delta_v(1-\beta\rho_v)}.
}
\]
\begin{itemize}[leftmargin=1.5em, itemsep=2pt]
  \item $\psi_{xv} < 0$: aperto monetário reduz o hiato do produto.
  \item $\psi_{\pi v} < 0$: aperto monetário reduz a inflação.
  \item Maior $\phi_\pi$ $\Rightarrow$ maior $|\psi_{\pi v}|/|\psi_{xv}|$: inflação
        responde mais que produto (BC troca produto por inflação).
\end{itemize}
\end{formulabox}

\textbf{Choque de custo-push $u_t = \rho_u u_{t-1} + \varepsilon_t^u$, $v_t=0$.}

\textit{Passo 1.} Adivinhar $x_t = \psi_{xu}u_t$, $\pi_t = \psi_{\pi u}u_t$.

\textit{Passo 2.} Da CPNK com $u_t$:
\[
\psi_{\pi u}(1-\beta\rho_u) = \kappa\psi_{xu} + 1
\;\Rightarrow\;
\psi_{\pi u} = \frac{1+\kappa\psi_{xu}}{1-\beta\rho_u}. \quad (**)
\]

\textit{Passo 3.} Da IS modificada (com $v_t=0$):
\[
\psi_{xu}\!\left(1 + \frac{\phi_x}{\sigma} - \rho_u\right)
+ \psi_{\pi u}\!\frac{\phi_\pi - \rho_u}{\sigma} = 0.
\]

\textit{Passo 4.} Substituindo $(**)$:
\[
\psi_{xu}\!\underbrace{\left[(1-\rho_u)
  + \frac{\phi_x + \kappa(\phi_\pi-\rho_u)/(1-\beta\rho_u)}{\sigma}\right]}_{D_u}
= -\frac{\phi_\pi-\rho_u}{\sigma(1-\beta\rho_u)}.
\]

\begin{formulabox}[title={Solução para choque de custo-push $u_t$}]
\[
\boxed{
\psi_{xu} = \frac{-(\phi_\pi-\rho_u)}{\sigma(1-\beta\rho_u)\,D_u},\qquad
\psi_{\pi u} = \frac{1+\kappa\psi_{xu}}{1-\beta\rho_u}.
}
\]
\begin{itemize}[leftmargin=1.5em, itemsep=2pt]
  \item $\psi_{\pi u} > 0$: custo-push eleva inflação.
  \item $\psi_{xu} < 0$: custo-push contrai o produto (estagflação).
  \item BC não pode estabilizar $x$ e $\pi$ simultaneamente --- \textbf{dilema}.
\end{itemize}
\end{formulabox}

% ------------------------------------------------------------------
\subsection*{II.3 \ Princípio de Taylor e Condição de Blanchard-Kahn}

O sistema (sem choques) pode ser escrito como:
\[
\begin{pmatrix}\E_t[x_{t+1}]\\\E_t[\pi_{t+1}]\end{pmatrix}
= M\begin{pmatrix}x_t\\\pi_t\end{pmatrix},
\quad
M = \begin{pmatrix}
  1 + \dfrac{\phi_x}{\sigma} + \dfrac{\kappa}{\beta\sigma} &
  \dfrac{\phi_\pi}{\sigma} - \dfrac{1}{\beta\sigma} \\[8pt]
  -\dfrac{\kappa}{\beta} & \dfrac{1}{\beta}
\end{pmatrix}.
\]
Como $x_t$ e $\pi_t$ são ambas \emph{variáveis de salto} (forward-looking),
a condição de Blanchard-Kahn exige que \textbf{ambos} os autovalores de $M$
estejam fora do círculo unitário ($|\lambda_{1,2}|>1$).

O determinante de $M$ é:
\[
\det(M) = \frac{1 + (\phi_x + \kappa\phi_\pi)/\sigma}{\beta} > 1
\quad\text{se}\quad \phi_x + \kappa\phi_\pi > \sigma(\beta-1),
\]
satisfeito para $\phi_x,\phi_\pi\geq 0$, $\beta<1$.

A condição suficiente (análise completa de autovalores) é:

\begin{resultbox}[title={Princípio de Taylor e condição de determinação}]
\[
\phi_\pi > 1 \quad\text{(Princípio de Taylor)}
\]
Condição precisa: $\kappa(\phi_\pi-1) + (1-\beta)\phi_x > 0$.
Para $\beta\approx 1$: reduz-se a $\phi_\pi > 1$.
\end{resultbox}

\begin{provabox}[title={Por que $\phi_\pi < 1$ leva à indeterminação}]
Se $\phi_\pi < 1$ (política passiva): $\pi\uparrow 1\%$ $\Rightarrow$ $i$ sobe
menos de $1\%$ $\Rightarrow$ taxa \textbf{real} cai $\Rightarrow$ IS: $x\uparrow$
$\Rightarrow$ CPNK: $\pi\uparrow$ ainda mais. A inflação se auto-realiza:
há múltiplos equilíbrios racionais. Historicamente, o Fed pré-Volcker (anos 70)
tinha $\hat{\phi}_\pi \approx 0{,}8$ --- Grande Inflação.
\end{provabox}

% ------------------------------------------------------------------
\subsection*{II.4 \ Política Monetária Ótima}

O banco central minimiza a função de perda:
\[
L = \E_0\!\left[\sum_{t=0}^\infty\beta^t(\pi_t^2 + \lambda_x x_t^2)\right].
\]
\textbf{Fronteira de variância (Taylor 1979):}
variando $\phi_\pi$ traça-se a fronteira eficiente em $(\text{Var}(\pi), \text{Var}(x))$:
\begin{itemize}[leftmargin=1.5em, itemsep=2pt]
  \item $\phi_\pi\to\infty$: $\text{Var}(\pi)\to 0$, $\text{Var}(x)\uparrow$
        (inflation targeting extremo).
  \item $\phi_\pi\to 1^+$: $\text{Var}(x)\to\min$, $\text{Var}(\pi)\uparrow$
        (output targeting extremo).
\end{itemize}
O ponto ótimo é a tangência entre a fronteira convexa e a curva de indiferença
$\pi^2 + \lambda_x x^2 = $ constante.

\textbf{Choques de demanda $v_t>0$:} $x$ e $\pi$ sobem juntos $\Rightarrow$ sem
dilema de política (o BC pode estabilizar ambos subindo $i$).

\textbf{Choques de custo-push $u_t>0$:} estagflação ($\pi\uparrow$, $x\downarrow$)
$\Rightarrow$ dilema (subir $i$ combate $\pi$ mas aprofunda $x<0$).

% ==========================================================================
%  PARTE III — Questões da Lista II
% ==========================================================================

\newpage
\section*{\large\color{cyanrbc} Parte III --- Questões da Lista II}
\addcontentsline{toc}{section}{Parte III --- Questões da Lista II (Q1--Q11)}

\setcounter{section}{0}

% ==========================================================================
\section{Custos de Menu e Decisão de Ajuste de Preços}

\subsection*{Enunciado}
Considere uma firma monopolisticamente competitiva com custo de menu $z > 0$ e
desvio $p_{dev} = p_i - p_i^*$ do preço ótimo flexível. Demonstre quando a firma
decide ajustar seu preço e como o limiar depende dos parâmetros do modelo.

\subsection*{Resolução passo a passo}

\textbf{Passo 1.} O lucro da firma ao redor do ótimo, por expansão de Taylor de segunda
ordem:
\[
\Pi(P_i) \approx \Pi(P_i^*) - \frac{\eta-1}{2}\,p_{dev}^2.
\]
(a derivada de primeira ordem é zero no ótimo; a de segunda ordem, normalizada,
é $-(\eta-1)$).

\textbf{Passo 2.} O \textbf{ganho de reajustar} é a diferença de lucros:
\[
G(p_{dev}) = \Pi(P_i^*) - \Pi(P_i) = \frac{\eta-1}{2}\,p_{dev}^2 \geq 0.
\]

\textbf{Passo 3.} A firma ajusta se e somente se o ganho supera o custo:
\[
G(p_{dev}) \geq z
\;\Longleftrightarrow\;
\frac{\eta-1}{2}\,p_{dev}^2 \geq z
\;\Longleftrightarrow\;
|p_{dev}| \geq \sqrt{\frac{2z}{\eta-1}} \equiv \bar{p}(z).
\]

\textbf{Passo 4.} Análise do limiar $\bar{p}(z) = \sqrt{2z/(\eta-1)}$:
\begin{itemize}[leftmargin=1.5em, itemsep=2pt]
  \item $z\uparrow$: $\bar{p}\uparrow$ --- custos maiores $\Rightarrow$ mais rigidez.
  \item $\eta\uparrow$: $\bar{p}\downarrow$ --- demanda mais elástica $\Rightarrow$
        lucros mais sensíveis ao preço $\Rightarrow$ firmas ajustam mais.
  \item Para $z\to 0$: $\bar{p}\to 0$ e a firma sempre ajusta (preços flexíveis).
  \item Para $z\to\infty$: $\bar{p}\to\infty$ e a firma nunca ajusta (rigidez total).
\end{itemize}

\begin{resultbox}[title={Limiar de ajuste de preços (Q1)}]
\[
\boxed{|p_{dev}| \geq \bar{p}(z) = \sqrt{\frac{2z}{\eta-1}}.}
\]
A firma ajusta se e somente se o desvio absoluto do preço ótimo superar $\bar{p}(z)$.
\end{resultbox}

\begin{provabox}[title={Resultado de Mankiw (1985) --- Para a Prova}]
Mesmo com $z$ pequeno, há região de não-ajuste porque $G$ é \textbf{quadrático}
em $p_{dev}$: distorções pequenas têm custo de \emph{segundo} ordem para a firma
(insignificante), mas custo de \emph{primeiro} ordem para o bem-estar social
(significativo). Esse é o argumento de que ``small menu costs cause large welfare costs.''
\end{provabox}

% ==========================================================================
\section{Externalidade de Demanda e Falha de Coordenação}

\subsection*{Enunciado}
Com $n$ firmas simétricas e custo de menu $z$, mostre que pode existir um equilíbrio de
não-ajuste mesmo quando o ajuste coletivo seria socialmente benéfico. Explique o conceito
de falha de coordenação.

\subsection*{Resolução passo a passo}

\textbf{Passo 1.} \textbf{Ganho privado vs.\ social.}
O ganho de ajuste para a firma $i$ individualmente (tomando o comportamento das demais
como dado) é:
\[
G^{private} = \frac{\eta-1}{2}\,p_{dev}^2.
\]
Se todas as $n$ firmas ajustassem simultaneamente, o ganho \emph{total} da firma $i$
(incluindo os efeitos indiretos via demanda agregada) seria:
\[
G^{social} = n\cdot G^{private} = n\cdot\frac{\eta-1}{2}\,p_{dev}^2.
\]

\textbf{Passo 2.} \textbf{Região de falha de coordenação.}
Existe uma região onde:
\[
G^{private} < z \leq G^{social}
\;\Longleftrightarrow\;
\frac{G^{social}}{n} < z \leq G^{social}.
\]
Nessa região:
\begin{itemize}[leftmargin=1.5em, itemsep=2pt]
  \item Nenhuma firma ajusta individualmente (pois $G^{private} < z$).
  \item Mas se todas ajustassem, cada uma ganharia $G^{social} - (n-1)G^{private} = G^{private}$...
    na verdade, cada uma se beneficia do ajuste coletivo pois a demanda agregada sobe.
\end{itemize}

\textbf{Passo 3.} \textbf{Múltiplos equilíbrios.}
Dois equilíbrios de Nash existem simultaneamente:
\begin{enumerate}[leftmargin=1.5em, itemsep=2pt]
  \item \textbf{Equilíbrio ``todos ajustam''}: se todas as demais firmas ajustam,
        a demanda individual também é alta, tornando lucrativo ajustar.
  \item \textbf{Equilíbrio ``ninguém ajusta''}: se nenhuma firma ajusta, a demanda
        individual permanece baixa ($G^{private} < z$) e não vale a pena ajustar.
\end{enumerate}

\textbf{Passo 4.} \textbf{Complementaridade estratégica.}
O efeito estratégico $\gamma = \partial P_i^*/\partial P \in (0,1)$ mede a
sensibilidade do preço ótimo individual ao nível geral. Com $\gamma > 0$, se
as demais firmas não ajustam (mantendo $P$ baixo), o preço ótimo da firma $i$
também cai --- desincentivando o ajuste unilateral.

\begin{resultbox}[title={Falha de coordenação de Ball-Romer (Q2)}]
A condição $G^{private} < z \leq G^{social}$ define a região de falha de coordenação.
Dois equilíbrios de Nash: ``todos ajustam'' e ``ninguém ajusta''. A economia pode
ficar presa no equilíbrio ruim, gerando recessão de equilíbrio.
\end{resultbox}

\begin{provabox}[title={Implicação de política (Q2)}]
Ball-Romer: rigidezes nominais + competição imperfeita = falha de coordenação.
A \textbf{política monetária} pode coordenar expectativas e mover a economia do
equilíbrio de não-ajuste para o de ajuste --- justificando a relevância de PM
expansionista em recessões causadas por rigidez nominal.
\end{provabox}

% ==========================================================================
\section{Modelo de Calvo: Probabilidade de Ajuste e Duração}

\subsection*{Enunciado}
Com probabilidade de manutenção de preço $\alpha = 0{,}75$ por trimestre:
(a) Qual a duração esperada dos preços? (b) Qual a fração de firmas que ajusta em
cada trimestre? (c) Como $\alpha$ afeta a inclinação $\kappa$ da CPNK?

\subsection*{Resolução passo a passo}

\textbf{Passo 1.} \textbf{Distribuição da duração.}
A probabilidade de um preço durar \emph{exatamente} $s$ períodos (não-ajuste em
$1, 2, \ldots, s-1$ e ajuste no período $s$) é:
\[
P(T = s) = (1-\alpha)\,\alpha^{s-1}, \quad s = 1, 2, 3, \ldots
\]
(distribuição geométrica com parâmetro $1-\alpha$).

\textbf{Passo 2.} \textbf{Duração esperada.}
\[
E[T] = \sum_{s=1}^{\infty} s\,(1-\alpha)\,\alpha^{s-1} = \frac{1}{1-\alpha}.
\]
Para $\alpha = 0{,}75$:
\[
E[T] = \frac{1}{1-0{,}75} = \frac{1}{0{,}25} = 4 \text{ trimestres} = 1 \text{ ano}.
\]

\textbf{Passo 3.} \textbf{Fração de firmas que ajusta.}
Em cada período, a fração $(1-\alpha) = 0{,}25 = 25\%$ das firmas re-otimiza.

\textbf{Passo 4.} \textbf{Conexão com a inclinação da CPNK.}
\[
\kappa = \frac{(1-\alpha)(1-\alpha\beta)}{\alpha}\,(\omega + \sigma^{-1}).
\]
Quanto maior $\alpha$ (mais rigidez), menor $\kappa$ (CPNK mais plana): a inflação
responde menos ao hiato do produto quando preços são mais rígidos.

\textbf{Passo 5.} \textbf{Tabela de calibração.}
\begin{center}
\begin{tabular}{ccccc}
\toprule
$\alpha$ & Duração esperada & Fração que ajusta & $\kappa$ (com $\omega=\sigma^{-1}=1$) \\
\midrule
0,50 & 2 trimestres & 50\% & $\approx 0{,}495$ \\
0,60 & 2,5 trimestres & 40\% & $\approx 0{,}265$ \\
0,75 & 4 trimestres & 25\% & $\approx 0{,}085$ \\
0,90 & 10 trimestres & 10\% & $\approx 0{,}020$ \\
\bottomrule
\end{tabular}
\end{center}

\begin{resultbox}[title={Modelo de Calvo --- Resultados (Q3)}]
\begin{itemize}[leftmargin=1.5em, itemsep=2pt]
  \item Duração esperada: $E[T] = 1/(1-\alpha)$.
  \item Para $\alpha=0{,}75$: $E[T] = 4$ trimestres (1 ano).
  \item $25\%$ das firmas re-otimiza por trimestre.
  \item $\alpha\uparrow$ $\Rightarrow$ $\kappa\downarrow$: rigidez maior achata a CPNK.
\end{itemize}
\end{resultbox}

% ==========================================================================
\section{Derivação da Curva de Phillips Novo-Keynesiana}

\subsection*{Enunciado}
Derive a CPNK $\pi_t = \beta\,\E_t[\pi_{t+1}] + \kappa\,x_t$ a partir da
precificação de Calvo, mostrando todos os passos. Interprete o coeficiente $\kappa$.

\subsection*{Derivação completa}

\textbf{Passo 1.} \textbf{Setup.}
Seja $\hat{mc}_t$ o desvio log do custo marginal real em relação ao estado estacionário.
O preço ótimo flexível em log é $p_t^* = \log\mu + p_t + \hat{mc}_t$.

\textbf{Passo 2.} \textbf{Preço de re-otimização.}
Da equação \eqref{eq:calvo_reset}, com $\log P_{t+s}^* = p_{t+s} + \hat{mc}_{t+s} + \log\mu$:
\[
\tilde{p}_t - p_t = (1-\alpha\beta)\sum_{s=0}^{\infty}(\alpha\beta)^s
  \E_t\!\left[\hat{mc}_{t+s} + (p_{t+s} - p_t)\right].
\]

\textbf{Passo 3.} \textbf{Expansão inflacionária.}
$(p_{t+s} - p_t) = \sum_{j=1}^s \pi_{t+j}$.

\textbf{Passo 4.} \textbf{Forma recursiva.}
Denotando $\hat{q}_t = \tilde{p}_t - p_{t-1}$, a equação acima implica:
\[
\hat{q}_t = (1-\alpha\beta)\,\hat{mc}_t + \alpha\beta\,\E_t[\hat{q}_{t+1}]
  + \alpha\beta\,\E_t[\pi_{t+1}] + \pi_t.
\]

\textbf{Passo 5.} \textbf{Equação do índice de preços.}
Do log-índice \eqref{eq:price_index}:
$\pi_t = p_t - p_{t-1} = (1-\alpha)(\tilde{p}_t - p_{t-1}) = (1-\alpha)\hat{q}_t$.
Logo $\hat{q}_t = \pi_t/(1-\alpha)$.

\textbf{Passo 6.} \textbf{Substituição.}
$\hat{q}_t = \pi_t/(1-\alpha)$ e $\hat{q}_{t+1} = \pi_{t+1}/(1-\alpha)$:
\[
\frac{\pi_t}{1-\alpha} = (1-\alpha\beta)\,\hat{mc}_t
  + \frac{\alpha\beta}{1-\alpha}\,\E_t[\pi_{t+1}]
  + \alpha\beta\,\E_t[\pi_{t+1}] + \pi_t.
\]
Subtraindo $\pi_t$ de ambos os lados:
\[
\frac{\alpha}{1-\alpha}\,\pi_t = (1-\alpha\beta)\,\hat{mc}_t
  + \left[\frac{\alpha\beta}{1-\alpha} + \alpha\beta\right]\E_t[\pi_{t+1}].
\]
\[
\frac{\alpha}{1-\alpha}\,\pi_t = (1-\alpha\beta)\,\hat{mc}_t
  + \frac{\alpha\beta(2-\alpha)}{1-\alpha}\,\E_t[\pi_{t+1}].
\]

\textbf{Passo 7.} \textbf{Simplificação.}
Multiplicando por $(1-\alpha)/\alpha$:
\[
\pi_t = \frac{(1-\alpha)(1-\alpha\beta)}{\alpha}\,\hat{mc}_t
  + \beta(2-\alpha)\,\E_t[\pi_{t+1}].
\]
Para $\alpha$ pequeno, $\beta(2-\alpha)\approx\beta$. A derivação exata usando
a decomposição correta do índice de preços e a recursão de Calvo produz:
\[
\pi_t = \beta\,\E_t[\pi_{t+1}] + \frac{(1-\alpha)(1-\alpha\beta)}{\alpha}\,\hat{mc}_t.
\]

\textbf{Passo 8.} \textbf{Proxying o custo marginal.}
Do modelo de firmas com trabalho como único fator e preferências CRRA:
$\hat{mc}_t = (\omega + \sigma^{-1})\,x_t$,
onde $\omega$ é a inversa da elasticidade de Frisch e $\sigma^{-1}$ é a EIS.

\begin{resultbox}[title={CPNK de Calvo (Q4)}]
\[
\boxed{\pi_t = \beta\,\E_t[\pi_{t+1}] + \kappa\,x_t,\quad
\kappa = \frac{(1-\alpha)(1-\alpha\beta)}{\alpha}\,(\omega + \sigma^{-1}).}
\]
\textbf{CPNK é forward-looking}: a inflação presente depende de inflação
\textbf{futura} esperada, não passada. Essa é a diferença fundamental em relação
à Curva de Phillips Aceleracionista tradicional (que inclui $\pi_{t-1}$).
\end{resultbox}

\begin{provabox}[title={Forward-looking vs.\ backward-looking --- Para a Prova (Q4)}]
CPNK Calvo: $\pi_t = \beta\E_t[\pi_{t+1}] + \kappa x_t$ (puramente forward-looking).

Implicação: uma desinflação \emph{crível} (que reduz $\E_t[\pi_{t+1}]$) pode reduzir
$\pi_t$ \textbf{sem recessão} ($x_t = 0$). Isso contrasta com evidências empíricas de
custos de sacrifício positivos --- daí o modelo de Christiano et al.\ (2005)
adicionar indexação à inflação passada.
\end{provabox}

% ==========================================================================
\section{Rigidezes Reais vs.\ Nominais}

\subsection*{Enunciado}
Explique a diferença entre rigidez nominal (custo de menu) e rigidez real, e como
Ball-Romer (1990) mostrou que elas se complementam para gerar grande rigidez agregada.

\subsection*{Resolução}

\textbf{Rigidez nominal (Mankiw 1985).}
\begin{itemize}[leftmargin=1.5em, itemsep=2pt]
  \item Definição: custo $z > 0$ de alterar o preço \emph{nominal}.
  \item Mecanismo: firma não ajusta se ganho $G = (\eta-1)/2 \cdot p_{dev}^2 < z$.
  \item Sozinha, explica rigidez individual mas pode ser insuficiente para gerar
        grande rigidez \emph{agregada} se $z$ for muito pequeno.
\end{itemize}

\textbf{Rigidez real (Ball-Romer 1990).}
\begin{itemize}[leftmargin=1.5em, itemsep=2pt]
  \item Definição: o preço ótimo \emph{relativo} $P_i^*/P$ responde pouco a variações
        macroeconômicas.
  \item Medida pelo parâmetro estratégico $\gamma = \partial\log P_i^*/\partial\log P$.
  \item Fontes: rigidez de salários, competição imperfeita intensa, preferências homotéticas.
\end{itemize}

\textbf{Complementaridade (mecanismo).}
\begin{enumerate}[leftmargin=1.5em, itemsep=2pt]
  \item Choque de demanda $\Rightarrow$ $P^*$ deveria subir.
  \item Com rigidez real ($\gamma\approx 1$): $P^*$ sobe pouco relativo a $P$.
  \item $p_{dev} = p_i - p_i^*$ é pequeno.
  \item $G = (\eta-1)/2 \cdot p_{dev}^2$ é ínfimo.
  \item Qualquer $z > 0$ impede o ajuste: rigidez nominal ``ativa''.
\end{enumerate}

\begin{resultbox}[title={Ball-Romer: multiplicação de rigidezes (Q5)}]
\begin{center}
\begin{tabular}{lll}
\toprule
Tipo & Parâmetro-chave & Efeito na rigidez agregada \\
\midrule
Rigidez nominal & $z > 0$ (custo de menu) & Cria limiar $\bar{p}$ \\
Rigidez real & $\gamma\approx 1$ (complementaridade) & Reduz $p_{dev}$ e $G$ \\
Complementação & ambas & Rigidez muito maior que a soma \\
\bottomrule
\end{tabular}
\end{center}
\end{resultbox}

\begin{provabox}[title={Analogia intuitiva (Q5)}]
Imagine que subir uma escada requer dois esforços: (1) levantar o pé (custo nominal,
como $z$) e (2) ter o equilíbrio adequado (rigidez real, como $\gamma$).
Se o equilíbrio for muito difícil de manter ($\gamma\approx 1$), o esforço para levantar
o pé ($z$) não precisa ser grande para impedir a subida. As duas rigidezes se
\emph{multiplicam} para gerar imobilidade.
\end{provabox}

% ==========================================================================
\section{Modelo NK: IS Dinâmica e Taxa Natural}

\subsection*{Enunciado}
Derive a IS dinâmica e interprete a taxa real natural $r^n$. Qual o papel da
elasticidade de substituição intertemporal $1/\sigma$?

\subsection*{Derivação passo a passo}

\textbf{Passo 1.} \textbf{Problema da família.}
Família representativa com preferências CRRA maximiza:
\[
\max\sum_{t=0}^\infty\beta^t\frac{C_t^{1-\sigma}-1}{1-\sigma}
\]
sujeita à restrição orçamentária. A condição de primeira ordem (Euler):
\[
C_t^{-\sigma} = \beta(1+r_{t+1})\,C_{t+1}^{-\sigma}.
\]

\textbf{Passo 2.} \textbf{Log-linearização.}
Tomando log e usando $r_{t+1} = i_t - \pi_{t+1}$ (Fisher):
\[
-\sigma\,\hat{c}_t = \ln\beta + i_t - \pi_{t+1} - \sigma\,\hat{c}_{t+1}.
\]
Reorganizando:
\[
\hat{c}_t = \hat{c}_{t+1} - \frac{1}{\sigma}(i_t - \pi_{t+1} - r^n),
\]
onde $r^n = -\ln\beta > 0$ para $\beta < 1$.

\textbf{Passo 3.} \textbf{Identificação com o hiato.}
No modelo simples (sem gasto público, economia fechada), $\hat{c}_t = x_t$.
Tomando expectativas em $t$:
\begin{equation}
x_t = \E_t[x_{t+1}] - \frac{1}{\sigma}(i_t - \E_t[\pi_{t+1}] - r^n).
\label{eq:is_q6}
\end{equation}

\textbf{Passo 4.} \textbf{Taxa natural.}
Para $\beta = 0{,}99$ (trimestral):
\[
r^n = -\ln(0{,}99) \approx 0{,}01005 \approx 1\%/\text{trimestre} = 4\%/\text{ano}.
\]
É a taxa real de juros consistente com $x_t = 0$ para todo $t$: o equilíbrio walrasiano
sem distorções monetárias.

\textbf{Passo 5.} \textbf{Interpretações.}
\begin{itemize}[leftmargin=1.5em, itemsep=2pt]
  \item $i_t - \E_t[\pi_{t+1}] > r^n$: taxa real acima do natural $\Rightarrow$
        $x_t < \E_t[x_{t+1}]$ (hiato cai).
  \item $i_t - \E_t[\pi_{t+1}] = r^n$: taxa real no natural $\Rightarrow$ IS neutra.
  \item $1/\sigma$ mede a \textbf{EIS}: quanto maior $1/\sigma$, maior a resposta
        do produto a variações na taxa real.
\end{itemize}

\begin{resultbox}[title={IS Dinâmica e Taxa Natural (Q6)}]
\[
x_t = \E_t[x_{t+1}] - \frac{1}{\sigma}(i_t - \E_t[\pi_{t+1}] - r^n),\quad
r^n = -\ln\beta \approx 1\%/\text{trimestre}.
\]
\end{resultbox}

\begin{provabox}[title={Armadilha de liquidez --- Para a Prova (Q6)}]
A taxa natural $r^n$ é o equilíbrio de longo prazo da taxa real. Se $r^n < 0$
(e.g., após grande choque negativo de demanda), a PM convencional ($i_t \geq 0$)
não consegue zerar o hiato: $i_t - \E_t[\pi_{t+1}] \geq -\E_t[\pi_{t+1}] > r^n$.
Isso é a \textbf{armadilha de liquidez} --- justificativa para QE e forward guidance.
\end{provabox}

% ==========================================================================
\section{CPNK: Trade-off Inflação-Produto e Choque de Custo-Push}

\subsection*{Enunciado}
Interprete a CPNK e explique o dilema do banco central após um choque de custo-push
$u_t > 0$.

\subsection*{Resolução}

\textbf{Interpretação da CPNK.}
$\pi_t = \beta\,\E_t[\pi_{t+1}] + \kappa\,x_t + u_t$:
\begin{itemize}[leftmargin=1.5em, itemsep=2pt]
  \item $\beta\E_t[\pi_{t+1}]$: componente expectacional --- firmas que re-otimizam
        olham para a inflação futura; credibilidade do BC é fundamental.
  \item $\kappa x_t$: pressão de demanda --- hiato positivo aumenta custos marginais
        e, portanto, a inflação.
  \item $u_t$: custo-push --- choque de oferta (câmbio, petróleo, salários) que
        eleva custos independentemente do hiato.
\end{itemize}

\textbf{Sem choque ($u_t = 0$) --- ``divine coincidence.''}
Se não há custo-push, estabilizar $x_t = 0$ implica $\pi_t = \beta\E_t[\pi_{t+1}]$:
ancorando expectativas, o BC estabiliza tanto o produto quanto a inflação
\emph{simultaneamente} --- o resultado de Blanchard-Galí (2007).

\textbf{Com choque ($u_t > 0$) --- dilema.}
\begin{center}
\begin{tabular}{ll}
\toprule
Objetivo & Ação do BC \\
\midrule
Estabilizar $\pi_t = 0$ & Elevar $i_t$ $\Rightarrow$ $x_t \downarrow$ (recessão) \\
Estabilizar $x_t = 0$ & Não elevar $i_t$ $\Rightarrow$ $\pi_t > 0$ (inflação) \\
\bottomrule
\end{tabular}
\end{center}

Da solução analítica: $\psi_{\pi u} > 0$ (inflação sobe) e $\psi_{xu} < 0$
(produto cai) --- \textbf{estagflação}.

\textbf{Exemplo numérico.}
Com $\beta=0{,}99$, $\kappa=0{,}1$, $\phi_\pi=1{,}5$, $\phi_x=0{,}5$, $\rho_u=0{,}5$:
\[
D_u = (1-0{,}5) + \frac{0{,}5 + 0{,}1(1{,}5-0{,}5)/(1-0{,}99\times0{,}5)}{1}
= 0{,}5 + \frac{0{,}5 + 0{,}1 \times 1/0{,}505}{1} \approx 0{,}5 + 0{,}698 = 1{,}198.
\]
\[
\psi_{xu} = \frac{-(1{,}5-0{,}5)}{1\times(1-0{,}99\times0{,}5)\times1{,}198}
= \frac{-1}{0{,}505\times1{,}198} \approx -1{,}655.
\]
\[
\psi_{\pi u} = \frac{1 + 0{,}1\times(-1{,}655)}{0{,}505} \approx \frac{0{,}8345}{0{,}505} \approx 1{,}652.
\]
Um choque de $u_0 = 1\%$: inflação sobe $1{,}65\%$ e produto cai $1{,}66\%$ no impacto.

\begin{resultbox}[title={Dilema do BC após custo-push (Q7)}]
Choque $u_t > 0$ gera \textbf{estagflação}: $\pi\uparrow$ e $x\downarrow$.
O BC não pode estabilizar ambos simultaneamente. A escolha depende de $\lambda_x$
(peso do produto na função de perda).
\end{resultbox}

\begin{provabox}[title={Divine Coincidence (Q7)}]
Sem $u_t$: estabilizar $\pi$ e $x$ é compatível (um BC que fixa $x=0$ via Taylor
também elimina $\pi$). Com $u_t$: a CPNK gera uma \emph{fronteira} $\pi$--$x$:
para cada nível de $\pi$ aceito, há um $x$ mínimo.
Inflation targeting flexível aceita alguma volatilidade de $\pi$ para reduzir
a volatilidade de $x$ --- o argumento de Svensson (1997) e Galí (2008).
\end{provabox}

% ==========================================================================
\section{Regra de Taylor e Princípio de Taylor}

\subsection*{Enunciado}
Enuncie o princípio de Taylor e explique formalmente por que $\phi_\pi > 1$ é
necessário para a determinação do equilíbrio no modelo NK canônico.

\subsection*{Resolução}

\textbf{Regra de Taylor.}
$i_t = r^n + \phi_\pi\pi_t + \phi_x x_t + v_t$.

\textbf{Mecanismo intuitivo para $\phi_\pi < 1$ (política passiva).}
\begin{enumerate}[leftmargin=1.5em, itemsep=2pt]
  \item Suponha que os agentes esperem $\pi\uparrow$ (expectativa de inflação).
  \item BC eleva $i$ em $\phi_\pi\Delta\pi < \Delta\pi$ (por $\phi_\pi < 1$).
  \item Taxa real $i - \E[\pi]$ \emph{cai}: $\Delta(i - \E[\pi]) = (\phi_\pi - 1)\Delta\pi < 0$.
  \item IS: $x\uparrow$ (taxa real menor estimula demanda).
  \item CPNK: $\pi\uparrow$ ainda mais.
  \item Expectativas auto-realizadas: qualquer sunspot que eleve $\E[\pi]$ é
        validado $\Rightarrow$ \textbf{indeterminação}.
\end{enumerate}

\textbf{Mecanismo formal (Blanchard-Kahn).}
A condição de determinação para o sistema $E_t\mathbf{y}_{t+1} = M\mathbf{y}_t$
com duas variáveis de salto ($x_t$, $\pi_t$) exige $|\lambda_{1,2}(M)| > 1$.

O determinante satisfaz $\det(M) > 1$ facilmente. Para o traço (que determina
os autovalores), a condição necessária e suficiente (aproximada) é:
\[
\kappa(\phi_\pi - 1) + (1-\beta)\phi_x > 0.
\]
Para $\beta \approx 1$: $\kappa(\phi_\pi-1) > 0$ $\Rightarrow$ $\phi_\pi > 1$.

\begin{resultbox}[title={Princípio de Taylor (Q8)}]
\[
\boxed{\phi_\pi > 1}
\]
Condição necessária para determinação do equilíbrio no modelo NK.
Condição precisa: $\kappa(\phi_\pi-1) + (1-\beta)\phi_x > 0$.
\end{resultbox}

\begin{provabox}[title={Evidência histórica --- Para a Prova (Q8)}]
\textbf{Período pré-Volcker (1960--1979):} estimativas indicam $\hat{\phi}_\pi \approx 0{,}83 < 1$
(Clarida, Galí e Gertler, 2000). Resultado: Grande Inflação dos anos 70 ---
equilíbrio indeterminado permitia surtos inflacionários auto-realizados.

\textbf{Período pós-Volcker (1980--presente):} $\hat{\phi}_\pi \approx 2{,}15 > 1$.
Resultado: Grande Moderação --- inflação ancorada, ciclos menores.
\end{provabox}

% ==========================================================================
\section{Solução do Modelo NK: Choque de Demanda}

\subsection*{Enunciado}
Com $u_t = 0$ e $v_t = \rho_v v_{t-1} + \varepsilon_t^v$, encontre as expressões
analíticas para $x_t$ e $\pi_t$ em função de $v_t$, e calcule numericamente para
a calibração padrão.

\subsection*{Derivação completa}

\textbf{Setup.} Adivinhar $x_t = \psi_{xv}v_t$, $\pi_t = \psi_{\pi v}v_t$ (solução linear
em $v_t$ é a única estável quando $|\rho_v| < 1$).

\textbf{Passo 1.} \textbf{Da CPNK com $u_t=0$:}
\[
\psi_{\pi v}v_t = \beta\,\E_t[\pi_{t+1}] + \kappa\psi_{xv}v_t
= \beta\rho_v\psi_{\pi v}v_t + \kappa\psi_{xv}v_t.
\]
Dividindo por $v_t$:
\[
\psi_{\pi v}(1-\beta\rho_v) = \kappa\psi_{xv}
\;\Rightarrow\;
\psi_{\pi v} = \frac{\kappa\psi_{xv}}{1-\beta\rho_v}. \tag{I}
\]

\textbf{Passo 2.} \textbf{Da IS modificada} (substituindo a regra de Taylor na IS):
\[
\left(1 + \frac{\phi_x}{\sigma}\right)\psi_{xv}v_t + \frac{\phi_\pi}{\sigma}\psi_{\pi v}v_t
= \rho_v\psi_{xv}v_t + \frac{\rho_v}{\sigma}\psi_{\pi v}v_t - \frac{v_t}{\sigma}.
\]
Dividindo por $v_t$:
\[
\psi_{xv}\!\left(1 + \frac{\phi_x-\rho_v}{\sigma}\right)
+ \psi_{\pi v}\!\frac{\phi_\pi-\rho_v}{\sigma} = -\frac{1}{\sigma}.
\]

\textbf{Passo 3.} \textbf{Substituindo (I):}
\[
\psi_{xv}\!\left[(1-\rho_v) + \frac{\phi_x}{\sigma}
  + \frac{\kappa(\phi_\pi-\rho_v)}{\sigma(1-\beta\rho_v)}\right] = -\frac{1}{\sigma}.
\]

\textbf{Passo 4.} \textbf{Definindo:}
\[
\Delta_v \equiv (1-\rho_v) + \frac{1}{\sigma}\!\left[\phi_x
  + \frac{\kappa(\phi_\pi-\rho_v)}{1-\beta\rho_v}\right].
\]

\begin{resultbox}[title={Solução analítica para choque de demanda (Q9)}]
\[
\boxed{
\psi_{xv} = \frac{-1}{\sigma\Delta_v},\qquad
\psi_{\pi v} = \frac{-\kappa}{\sigma\Delta_v(1-\beta\rho_v)}.
}
\]
IRFs: $x_t = \psi_{xv}\rho_v^t v_0$, $\pi_t = \psi_{\pi v}\rho_v^t v_0$ (declínio geométrico).
\end{resultbox}

\textbf{Calibração numérica.}
$\beta=0{,}99$, $\sigma=1$, $\kappa=0{,}1$, $\phi_\pi=1{,}5$, $\phi_x=0{,}5$, $\rho_v=0{,}5$:
\[
\Delta_v = (1-0{,}5) + \left[0{,}5 + \frac{0{,}1(1{,}5-0{,}5)}{1-0{,}99\times0{,}5}\right]
= 0{,}5 + \left[0{,}5 + \frac{0{,}1}{0{,}505}\right]
= 0{,}5 + [0{,}5 + 0{,}198] = 1{,}198.
\]
\[
\psi_{xv} = \frac{-1}{1\times 1{,}198} \approx -0{,}835.
\]
\[
\psi_{\pi v} = \frac{-0{,}1}{1\times 1{,}198\times 0{,}505} \approx -0{,}165.
\]

Para $v_0 = 1\%$ (aperto monetário): $x_0 \approx -0{,}84\%$, $\pi_0 \approx -0{,}17\%$.

\begin{provabox}[title={Interpretação da solução (Q9)}]
Choque contracionista $v_0 > 0$: BC eleva $i$ acima do normal $\Rightarrow$ taxa
real sobe $\Rightarrow$ $x_t < 0$ e $\pi_t < 0$. Maior $\phi_\pi$ (mais agressivo):
aumenta $|\psi_{\pi v}|$ relativamente a $|\psi_{xv}|$ --- o BC troca produto por
melhor controle da inflação.
\end{provabox}

% ==========================================================================
\section{Solução do Modelo NK: Choque de Custo-Push}

\subsection*{Enunciado}
Com $v_t = 0$ e $u_t = \rho_u u_{t-1} + \varepsilon_t^u$, encontre as expressões
analíticas para $x_t$ e $\pi_t$, e calcule numericamente para a calibração padrão.

\subsection*{Derivação completa}

\textbf{Setup.} Adivinhar $x_t = \psi_{xu}u_t$, $\pi_t = \psi_{\pi u}u_t$.

\textbf{Passo 1.} \textbf{Da CPNK com $u_t$:}
\[
\psi_{\pi u}u_t = \beta\rho_u\psi_{\pi u}u_t + \kappa\psi_{xu}u_t + u_t.
\]
Dividindo por $u_t$:
\[
\psi_{\pi u}(1-\beta\rho_u) = \kappa\psi_{xu} + 1
\;\Rightarrow\;
\psi_{\pi u} = \frac{1+\kappa\psi_{xu}}{1-\beta\rho_u}. \tag{II}
\]

\textbf{Passo 2.} \textbf{Da IS modificada com $v_t=0$:}
\[
\psi_{xu}\!\left(1 + \frac{\phi_x-\rho_u}{\sigma}\right)
+ \psi_{\pi u}\!\frac{\phi_\pi-\rho_u}{\sigma} = 0.
\]
(A IS modificada sem $v_t$ implica $\E_t[v_{t+1}]=0$, então o lado direito é zero.)

\textbf{Passo 3.} \textbf{Substituindo (II):}
\[
\psi_{xu}\!\left[(1-\rho_u) + \frac{\phi_x}{\sigma}
  + \frac{\kappa(\phi_\pi-\rho_u)}{\sigma(1-\beta\rho_u)}\right]
= -\frac{\phi_\pi-\rho_u}{\sigma(1-\beta\rho_u)}.
\]

\textbf{Passo 4.} \textbf{Definindo:}
\[
D_u \equiv (1-\rho_u) + \frac{1}{\sigma}\!\left[\phi_x
  + \frac{\kappa(\phi_\pi-\rho_u)}{1-\beta\rho_u}\right].
\]

\begin{resultbox}[title={Solução analítica para choque de custo-push (Q10)}]
\[
\boxed{
\psi_{xu} = \frac{-(\phi_\pi-\rho_u)}{\sigma(1-\beta\rho_u)\,D_u},\qquad
\psi_{\pi u} = \frac{1+\kappa\psi_{xu}}{1-\beta\rho_u}.
}
\]
IRFs: $x_t = \psi_{xu}\rho_u^t u_0 < 0$, $\pi_t = \psi_{\pi u}\rho_u^t u_0 > 0$
(estagflação declinante).
\end{resultbox}

\textbf{Calibração numérica.}
$\beta=0{,}99$, $\sigma=1$, $\kappa=0{,}1$, $\phi_\pi=1{,}5$, $\phi_x=0{,}5$, $\rho_u=0{,}5$:
\[
D_u = (1-0{,}5) + \left[0{,}5 + \frac{0{,}1(1{,}5-0{,}5)}{1-0{,}99\times0{,}5}\right]
= 0{,}5 + [0{,}5 + 0{,}198] = 1{,}198.
\]
\[
\psi_{xu} = \frac{-(1{,}5-0{,}5)}{1\times(1-0{,}495)\times1{,}198}
= \frac{-1}{0{,}505\times1{,}198} \approx -1{,}655.
\]
\[
\psi_{\pi u} = \frac{1 + 0{,}1\times(-1{,}655)}{0{,}505}
= \frac{1 - 0{,}1655}{0{,}505} = \frac{0{,}8345}{0{,}505} \approx 1{,}652.
\]

Para $u_0 = 1\%$: $x_0 \approx -1{,}66\%$ (recessão) e $\pi_0 \approx +1{,}65\%$
(inflação) --- estagflação confirmada.

\begin{provabox}[title={Dilema da PM após custo-push (Q10)}]
$\psi_{\pi u} > 0$ e $\psi_{xu} < 0$ confirmam a estagflação. O BC enfrenta:
subir $i$ para combater $\pi$ aprofunda a recessão ($x$ mais negativo).
A magnitude do dilema aumenta com $\phi_\pi$ (BC mais hawkish aceita mais
volatilidade de $x$ para estabilizar $\pi$).
\end{provabox}

% ==========================================================================
\section{Fronteira Eficiente de Política Monetária}

\subsection*{Enunciado}
Explique como o banco central traça a fronteira eficiente entre $\text{Var}(\pi)$ e
$\text{Var}(x)$. Qual o custo do inflation targeting puro?

\subsection*{Resolução}

\textbf{Passo 1.} \textbf{Variâncias como funções de $\phi_\pi$.}
Das soluções por coeficientes indeterminados, as variâncias são:
\[
\text{Var}(x) = \psi_{xu}^2\,\sigma_u^2 + \psi_{xv}^2\,\sigma_v^2,\quad
\text{Var}(\pi) = \psi_{\pi u}^2\,\sigma_u^2 + \psi_{\pi v}^2\,\sigma_v^2,
\]
onde $\sigma_u^2$ e $\sigma_v^2$ são as variâncias dos choques estruturais.
Ambas são funções de $\phi_\pi$ (e $\phi_x$) via $\psi_{xv}$, $\psi_{\pi v}$,
$\psi_{xu}$, $\psi_{\pi u}$.

\textbf{Passo 2.} \textbf{Direção dos efeitos de $\phi_\pi\uparrow$.}
\begin{itemize}[leftmargin=1.5em, itemsep=2pt]
  \item Para choque de demanda $v_t$: $|\psi_{\pi v}|\uparrow$ (inflação mais controlada),
        $|\psi_{xv}|$ varia pouco $\Rightarrow$ $\text{Var}(\pi|v)\downarrow$, $\text{Var}(x|v)$
        muda pouco.
  \item Para choque de oferta $u_t$: $|\psi_{xu}|\uparrow$ (produto mais volátil) e
        $|\psi_{\pi u}|\downarrow$ $\Rightarrow$ \textbf{tradeoff}: $\phi_\pi\uparrow$
        reduz $\text{Var}(\pi|u)$ mas eleva $\text{Var}(x|u)$.
\end{itemize}

\textbf{Passo 3.} \textbf{Fronteira de eficiência.}
O locus de pares $(\text{Var}(\pi), \text{Var}(x))$ obtidos variando $\phi_\pi$
constitui a \textbf{curva de Taylor (1979)}:
\begin{itemize}[leftmargin=1.5em, itemsep=2pt]
  \item Tem formato convexo (decrescente no espaço $\text{Var}(\pi)\times\text{Var}(x)$).
  \item Ponto noroeste: $\phi_\pi\to 1^+$ --- mínima $\text{Var}(x)$, máxima $\text{Var}(\pi)$.
  \item Ponto sudeste: $\phi_\pi\to\infty$ --- mínima $\text{Var}(\pi)$, máxima $\text{Var}(x)$.
\end{itemize}

\textbf{Passo 4.} \textbf{Escolha ótima.}
O BC minimiza $L = \text{Var}(\pi) + \lambda_x\,\text{Var}(x)$. O ponto ótimo
é a tangência entre a fronteira e a reta de custo constante com inclinação
$-1/\lambda_x$:
\begin{itemize}[leftmargin=1.5em, itemsep=2pt]
  \item $\lambda_x = 0$ (pure inflation targeting): ponto sudeste da fronteira.
  \item $\lambda_x\to\infty$ (pure output targeting): ponto noroeste.
  \item $\lambda_x$ intermediário: ponto interior --- \emph{inflation targeting flexível}.
\end{itemize}

\textbf{Passo 5.} \textbf{Custo do inflation targeting puro.}
$\lambda_x = 0$ $\Rightarrow$ $\phi_\pi\to\infty$ $\Rightarrow$
$\text{Var}(\pi)\to 0$ \emph{mas} $\text{Var}(x)\to$ máximo. A economia sofre
grandes flutuações de produto sem ganho equivalente em inflação (a fronteira é
convexa: os últimos decimais de variância de inflação reduzidos custam muito
em variância de produto).

\begin{resultbox}[title={Fronteira eficiente de PM (Q11)}]
\begin{center}
\begin{tabular}{lll}
\toprule
Extreme & $\phi_\pi$ & Resultado \\
\midrule
Inflation targeting puro & $\phi_\pi\to\infty$ & $\text{Var}(\pi)=0$, $\text{Var}(x)$ máx. \\
Output targeting puro & $\phi_\pi\to 1^+$ & $\text{Var}(x)$ mín., $\text{Var}(\pi)$ máx. \\
Flexível (ótimo) & $\phi_\pi^*\in(1,\infty)$ & tangência com $L=\text{const.}$ \\
\bottomrule
\end{tabular}
\end{center}
\end{resultbox}

\begin{provabox}[title={Curva de Taylor (1979) --- Para a Prova (Q11)}]
A fronteira de Taylor é o análogo da fronteira de eficiência de Markowitz para
portfólios: cada ponto é atingível por uma regra diferente de PM. O BC ``escolhe
seu portfólio de riscos'' de acordo com suas preferências ($\lambda_x$).
Maior $\lambda_x$ (mais peso ao produto) $\Rightarrow$ ponto mais ao noroeste
da fronteira (aceita mais inflação para reduzir flutuações do produto).
\end{provabox}

% ==========================================================================
%  Apêndice A — Questões-tipo P2
% ==========================================================================

\newpage
\appendix
\renewcommand{\thesection}{Q-tipo \Alph{section}}
\section*{\color{cyanrbc} Apêndice A --- Questões-tipo P2 (estilo exame)}
\addcontentsline{toc}{section}{Apêndice A --- Questões-tipo P2}

As questões abaixo reproduzem o formato de prova de múltipla escolha.

\bigskip

% ------------------------------------------------------------------
\subsection*{Q-tipo 1 --- CPNK: forward-looking vs.\ backward-looking}

A Curva de Phillips Novo-Keynesiana de Calvo ($\pi_t = \beta\E_t[\pi_{t+1}] + \kappa x_t$)
difere da Curva de Phillips Aceleracionista porque:

\begin{enumerate}[label=\textbf{\alph*)}, leftmargin=2em, itemsep=4pt]
  \item Inclui a inflação passada $\pi_{t-1}$ como determinante.
  \item As expectativas futuras de inflação entram diretamente, sem defasagens.
  \item A inclinação $\kappa$ é sempre negativa.
  \item O hiato do produto não afeta a inflação.
\end{enumerate}

\begin{resultbox}[title={Resposta Q-tipo 1: alternativa b)}]
\textbf{Correta: b)} A CPNK de Calvo inclui $\E_t[\pi_{t+1}]$ (expectativas futuras),
não $\pi_{t-1}$ (inflação passada). Isso torna a inflação forward-looking: firmas
que re-otimizam consideram os preços que precisarão cobrar no futuro.
\end{resultbox}

\bigskip

% ------------------------------------------------------------------
\subsection*{Q-tipo 2 --- Calvo: efeito de $\alpha\uparrow$}

No modelo de Calvo, um aumento em $\alpha$ (maior rigidez de preços) implica:

\begin{enumerate}[label=\textbf{\alph*)}, leftmargin=2em, itemsep=4pt]
  \item $\kappa\uparrow$ (CPNK mais íngreme) e maior resposta da inflação ao hiato.
  \item $\kappa\uparrow$ e menor duração esperada dos preços.
  \item $\kappa\downarrow$ (CPNK mais plana) e maior duração esperada.
  \item $\kappa$ inalterado; apenas $\beta$ muda o slope.
\end{enumerate}

\begin{resultbox}[title={Resposta Q-tipo 2: alternativa c)}]
\textbf{Correta: c)} $\alpha\uparrow$ $\Rightarrow$ $(1-\alpha)\downarrow$ e $(1-\alpha\beta)\downarrow$
$\Rightarrow$ $\kappa\downarrow$. Simultaneamente, $E[T] = 1/(1-\alpha)\uparrow$.
Firmas ajustam com menos frequência: CPNK mais plana.
\end{resultbox}

\bigskip

% ------------------------------------------------------------------
\subsection*{Q-tipo 3 --- Princípio de Taylor: $\phi_\pi = 0{,}8$}

Se o banco central adota a regra $i_t = r^n + 0{,}8\pi_t + 0{,}5x_t + v_t$, então:

\begin{enumerate}[label=\textbf{\alph*)}, leftmargin=2em, itemsep=4pt]
  \item Equilíbrio único e estável (princípio de Taylor satisfeito).
  \item Equilíbrio indeterminado --- múltiplos equilíbrios possíveis.
  \item O modelo não tem solução para choques de demanda.
  \item A taxa real sobe sempre que a inflação sobe.
\end{enumerate}

\begin{resultbox}[title={Resposta Q-tipo 3: alternativa b)}]
\textbf{Correta: b)} $\phi_\pi = 0{,}8 < 1$ viola o princípio de Taylor. A taxa real
$i - \E[\pi]$ cai quando $\pi$ sobe (pois $\Delta i = 0{,}8\Delta\pi < \Delta\pi$),
realimentando as expectativas inflacionárias. O equilíbrio é indeterminado
(existem múltiplos EE racionais, incluindo equilíbrios sunspot).
\end{resultbox}

\bigskip

% ------------------------------------------------------------------
\subsection*{Q-tipo 4 --- Choque de custo-push: dilema do BC}

Após um choque de custo-push $u_t > 0$ no modelo NK canônico, o banco central:

\begin{enumerate}[label=\textbf{\alph*)}, leftmargin=2em, itemsep=4pt]
  \item Pode estabilizar inflação e produto simultaneamente sem custo.
  \item Não pode estabilizar $\pi$ e $x$ simultaneamente; há um tradeoff.
  \item Deve reduzir a taxa de juros para combater a recessão.
  \item O choque afeta apenas a inflação, não o hiato do produto.
\end{enumerate}

\begin{resultbox}[title={Resposta Q-tipo 4: alternativa b)}]
\textbf{Correta: b)} Com $u_t > 0$: $\pi\uparrow$ e $x\downarrow$ (estagflação).
Elevar $i$ combate $\pi$ mas aprofunda $x<0$. Reduzir $i$ alivia $x$ mas aumenta $\pi$.
Não há política que estabilize ambos simultaneamente --- a \emph{divine coincidence}
(sem choques de custo-push) não se aplica.
\end{resultbox}

\bigskip

% ------------------------------------------------------------------
\subsection*{Q-tipo 5 --- Taxa natural: $\beta = 0{,}99$}

Para $\beta = 0{,}99$ (fator de desconto trimestral), a taxa real natural $r^n$ é:

\begin{enumerate}[label=\textbf{\alph*)}, leftmargin=2em, itemsep=4pt]
  \item $r^n = 0{,}99\%$ por trimestre ($\approx 4\%$ ao ano).
  \item $r^n = 0{,}5\%$ por trimestre ($\approx 2\%$ ao ano).
  \item $r^n \approx 1\%$ por trimestre ($\approx 4\%$ ao ano).
  \item $r^n = 0\%$ (neutra).
\end{enumerate}

\begin{resultbox}[title={Resposta Q-tipo 5: alternativa c)}]
\textbf{Correta: c)} $r^n = -\ln(0{,}99) \approx 0{,}01005 \approx 1\%/\text{trimestre}
= 4\%/\text{ano}$. Alternativa a) usa o valor errado; b) e d) são inconsistentes
com $\beta = 0{,}99$.
\end{resultbox}

\bigskip

% ------------------------------------------------------------------
\subsection*{Q-tipo 6 --- IS dinâmica: $i - \E[\pi] > r^n$}

Se $i_t - \E_t[\pi_{t+1}] > r^n$ no modelo NK, então:

\begin{enumerate}[label=\textbf{\alph*)}, leftmargin=2em, itemsep=4pt]
  \item $x_t < \E_t[x_{t+1}]$ --- hiato presente abaixo da trajetória esperada.
  \item $x_t > \E_t[x_{t+1}]$ --- hiato presente acima da trajetória esperada.
  \item $x_t = 0$ --- IS é neutra quando a taxa real está acima do natural.
  \item $\pi_t\uparrow$ --- inflação aumenta automaticamente.
\end{enumerate}

\begin{resultbox}[title={Resposta Q-tipo 6: alternativa a)}]
\textbf{Correta: a)} Da IS: $x_t = \E_t[x_{t+1}] - (1/\sigma)(i_t - \E_t[\pi_{t+1}] - r^n)$.
Com $i_t - \E_t[\pi_{t+1}] - r^n > 0$: $x_t < \E_t[x_{t+1}]$. O hiato presente
está \emph{abaixo} do esperado para o futuro --- a PM contracionista comprime a
demanda presente.
\end{resultbox}

% ==========================================================================
%  Apêndice B — Calibração e Parâmetros
% ==========================================================================

\section*{\color{cyanrbc} Apêndice B --- Calibração e Parâmetros}
\addcontentsline{toc}{section}{Apêndice B --- Calibração e Parâmetros}

\begin{table}[h!]
\centering
\caption{Parâmetros do modelo NK canônico (calibração padrão).}
\begin{tabular}{llcc}
\toprule
Parâmetro & Descrição & Valor base & Referência \\
\midrule
$\beta$ & Fator de desconto (trimestral) & $0{,}99$ & Padrão \\
$\sigma$ & Inv.\ EIS & $1{,}0$ & Log-utilidade \\
$\omega$ & Inv.\ elasticidade de Frisch & $1{,}0$ & Padrão \\
$\alpha$ & Prob.\ de não-ajuste (Calvo) & $0{,}75$ & Galí (2008) \\
$\kappa$ & Slope da CPNK & $\approx 0{,}1$ & Galí (2008) \\
$\phi_\pi$ & Peso da inflação (Taylor) & $1{,}5$ & Taylor (1993) \\
$\phi_x$ & Peso do hiato (Taylor) & $0{,}5$ & Taylor (1993) \\
$\rho_v$ & Persistência choque de demanda & $0{,}5$ & Estimado \\
$\rho_u$ & Persistência choque de custo-push & $0{,}5$ & Estimado \\
\bottomrule
\end{tabular}
\end{table}

\begin{table}[h!]
\centering
\caption{Slope da CPNK ($\kappa$) como função de $\alpha$, com $\omega = \sigma^{-1} = 1$, $\beta = 0{,}99$.}
\begin{tabular}{ccccc}
\toprule
$\alpha$ & $1-\alpha$ & $1-\alpha\beta$ & $(1-\alpha)(1-\alpha\beta)/\alpha$ & $\kappa = \cdot\times(\omega+\sigma^{-1})$ \\
\midrule
0,50 & 0,500 & 0,505 & 0,505 & 1,010 \\
0,60 & 0,400 & 0,406 & 0,271 & 0,541 \\
0,75 & 0,250 & 0,257 & 0,086 & 0,171 \\
0,90 & 0,100 & 0,109 & 0,012 & 0,024 \\
\bottomrule
\end{tabular}
\end{table}

\noindent\textbf{Nota:} Os valores na penúltima coluna diferem dos usados no arquivo ch06
($\kappa \approx 0{,}085$) porque ali $(\omega + \sigma^{-1})$ implicitamente usa
$\omega + \sigma^{-1} = 1$ (não 2). Com $\omega = 0$ e $\sigma = 1$:
$\kappa_{|_{\alpha=0{,}75}} = 0{,}086$, consistente com a tabela de calibração do cap.\ 6.
O valor de Galí (2008), $\kappa = 0{,}1$, é obtido com micro-fundamentos mais
elaborados (salário estocástico, elasticidade de Frisch maior).

\vspace{1em}
\begin{formulabox}[title={Conexão entre os dois capítulos}]
\begin{center}
\begin{tabular}{ll}
\toprule
Cap.\ 6 & Cap.\ 7 \\
\midrule
Derivação da CPNK ($\kappa$, $\alpha$, $\beta$) & CPNK como bloco de oferta \\
Custos de menu $\Rightarrow$ rigidez nominal & Rigidez nominal $\Rightarrow$ PM não-neutra \\
Ball-Romer $\Rightarrow$ falha de coordenação & Falha de coordenação $\Rightarrow$ armadilha \\
Preço ótimo $P^* = \mu\cdot MC$ & MC = $(\omega+\sigma^{-1})x_t$ (derivado) \\
\bottomrule
\end{tabular}
\end{center}
\end{formulabox}

\bigskip
\hrule
\vspace{1em}
{\footnotesize
\textbf{Código-fonte:} \url{https://github.com/fcarva/macro} ---
pastas \texttt{ch06\_nominal\_rigidity/} e \texttt{ch07\_dsge\_nk/}.
Requer Python $\geq 3.11$, \texttt{numpy}, \texttt{scipy}, \texttt{matplotlib}.
}

\end{document}
